This appendix is intentionally beyond the strict core of the booklet. It is included for students who want a clearer bridge to university mathematics, engineering, and Extension 2 style applications.

\subsubsection*{First-order linear equations}
At a higher level, many important models are written in the form
\[
\frac{dy}{dx} + P(x)y = Q(x).
\]
These are called first-order linear differential equations. They are usually solved by an integrating factor rather than by direct separation. This method is valuable in university mathematics, physics, and engineering, but it is not the central HSC focus here.

\subsubsection*{Stability language}
When an equilibrium solution is nearby, we can ask what happens to solutions that start close to it.

\begin{itemize}[leftmargin=*]
  \item An equilibrium is often called \textbf{stable} if nearby solutions move toward it.
  \item It is often called \textbf{unstable} if nearby solutions move away from it.
\end{itemize}

This language is common in differential equations, control theory, and dynamical systems.

\subsubsection*{Links to mechanics}
Extension 2 mechanics repeatedly uses rates of change. Velocity is the rate of change of position, and acceleration is the rate of change of velocity. As a result, many mechanics questions naturally become differential equations, especially when forces depend on displacement, velocity, or both.

\subsubsection*{Where this goes next}
Students moving further into mathematics, physics, economics, or engineering will later meet:
\begin{itemize}[leftmargin=*]
  \item linear differential equations,
  \item systems of differential equations,
  \item numerical methods,
  \item oscillation and stability,
  \item mathematical modelling of real systems.
\end{itemize}

The HSC treatment is only the beginning, but it introduces the most important habit early: connect algebra, graphs, and interpretation instead of treating them as separate skills.
